\begin{figure}[t]
\centering
\begin{tikzpicture}[x=1.35cm,y=1.05cm,
  phase/.style={rectangle,rounded corners,draw,fill=blue!8,font=\scriptsize,align=center,minimum width=2.2cm},
  risk/.style={font=\tiny,color=red!60!black,align=center}]
\node[phase] (write) at (0,0) {写入\\Write};
\node[phase] (store) at (2,0) {存储\\Store};
\node[phase] (retrieve) at (4,0) {检索\\Retrieve};
\node[phase] (use) at (6,0) {使用\\Use};
\node[phase] (forget) at (2,-1.4) {遗忘/回滚\\Forget};
\node[phase] (share) at (4,-1.4) {共享/传播\\Share};
\draw[->,>=stealth,thick] (write) -- (store);
\draw[->,>=stealth,thick] (store) -- (retrieve);
\draw[->,>=stealth,thick] (retrieve) -- (use);
\draw[->,>=stealth,thick] (store) -- (forget);
\draw[->,>=stealth,thick] (store) -- (share);
\draw[->,>=stealth,dashed,thick] (use) to[bend left=25] (write);
\node[risk] at (0,0.75) {错误写入\\投毒注入};
\node[risk] at (2,0.75) {摘要失真\\来源丢失};
\node[risk] at (4,0.75) {过时信息\\检索错配};
\node[risk] at (6,0.75) {错误使用\\过度泛化};
\node[risk] at (2,-2.05) {误删/残留衍生};
\node[risk] at (4,-2.05) {跨智能体污染};
\node[risk] at (3.2,1.45) {使用反馈驱动再写入（红色为本文关注的关键风险）};
\end{tikzpicture}
\caption{记忆生命周期与关键风险示意（概念图，依据综述文献的生命周期框架
归纳\cite{W0039}；红色标注各阶段风险，详见第~\ref{sec:update}~节）。}
\label{fig:lifecycle}
\end{figure}